
Language models are evaluated across multiple trustworthiness dimensions including truthfulness, fairness, and adversarial robustness. Standard evaluation practices measure these dimensions independently using separate benchmarks (TruthfulQA, BBQ, ANLI). However, interventions designed to improve performance on one dimension may affect others through shared neural representations. The nature and consistency of these cross-dimensional effects has not been systematically characterized.

This work investigates whether targeted interventions produce predictable cross-dimensional performance changes. The central question is whether dimension pairs exhibit consistent correlation patterns (trade-offs or co-improvement) that replicate across model architectures, or whether cross-dimensional effects are random and model-specific. Understanding these relationships is relevant for multi-objective trustworthiness optimization, where practitioners must predict unintended consequences of targeted improvements.

We conducted perturbation-based experiments applying LoRA fine-tuning to improve truthfulness, then measured resulting changes across three dimensions (truthfulness, fairness, robustness) and three model architectures (GPT-2, OPT, Pythia). The perturbation approach generates controlled variance by varying random seeds and hyperparameters, enabling correlation analysis that reveals dimension relationships. We also analyzed internal representation changes using Centered Kernel Alignment (CKA) to connect parameter updates to performance effects.

The study yields three main findings. First, LoRA interventions targeting attention layers cause universal representation changes (100\% of analyzed layers showed detectable changes). Second, cross-dimensional correlations vary by dimension pair: some pairs show near-zero correlation while others show strong negative correlation. Third, the direction and magnitude of correlations differ across model architectures, with partial replication of negative fairness-robustness correlation in two of three models tested.

This paper is organized as follows. Section 2 reviews related work on trustworthiness evaluation and multi-task learning. Section 3 describes the perturbation-based experimental methodology. Section 4 details experimental protocols for five hypotheses spanning existence, mechanism, and generalization. Section 5 presents results. Section 6 discusses interpretation and limitations. Section 7 concludes.
